\section{Theoretical Analysis}
\label{sec:theory}

We derive an upper bound on the forgetting of previously learned tasks under
sequential low-rank adaptation, and show that the FOLoRA regularizer minimizes its
dominant term.

\subsection{Setup}
Let $\theta_\tau$ denote the trainable parameters (adapter plus classifier) after
learning task $\mathcal{T}_\tau$, and let the adapter component be $\Delta W^{\tau}$.
For a later task $\mathcal{T}_t$ ($t>\tau$), let $\delta W=\Delta W^{t}-\Delta W^{\tau}$
be the adapter displacement. We measure the forgetting of task $\tau$ by the increase
of its empirical loss,
\begin{equation}
\Delta\mathcal{L}_\tau \;=\; \mathcal{L}_\tau(\theta_t) - \mathcal{L}_\tau(\theta_\tau).
\label{eq:forget_def}
\end{equation}

\subsection{A second-order forgetting bound}
\label{sec:bound}

We make the standard assumption (used by EWC~\cite{kirkpatrick2017overcoming} and
GPM~\cite{saha2021gradient}) that each task is trained to a (near-)stationary point, so
$\nabla_{\theta}\mathcal{L}_\tau(\theta_\tau)\approx 0$.

\begin{theorem}[Forgetting bound]
\label{thm:forgetting}
Suppose the empirical loss $\mathcal{L}_\tau$ is twice differentiable and $\theta_\tau$
is a stationary point of $\mathcal{L}_\tau$. Then, up to second order,
\begin{equation}
\Delta\mathcal{L}_\tau
\;\le\; \frac12\,\delta W^{\!\top} F_\tau\,\delta W + \mathcal{O}\big(\|\delta W\|^3\big),
\label{eq:bound1}
\end{equation}
where $F_\tau=\mathbb{E}_{(x,y)\sim\mathcal{D}_\tau}\big[\nabla\log p_\theta(y|x)\nabla\log p_\theta(y|x)^{\!\top}\big]$
is the empirical Fisher of task $\tau$ restricted to the adapter parameters.
\end{theorem}

\begin{proof}
Taylor-expanding $\mathcal{L}_\tau$ at $\theta_\tau$ along the displacement
$\delta\theta=\theta_t-\theta_\tau$ gives
\begin{equation}
\mathcal{L}_\tau(\theta_t)
= \mathcal{L}_\tau(\theta_\tau)
+ \langle\nabla\mathcal{L}_\tau(\theta_\tau),\delta\theta\rangle
+ \tfrac12\,\delta\theta^{\!\top} H_\tau\,\delta\theta
+ \mathcal{O}(\|\delta\theta\|^3),
\end{equation}
where $H_\tau$ is the Hessian. Since $\nabla\mathcal{L}_\tau(\theta_\tau)\approx0$, the
linear term vanishes. Under the Fisher-information approximation $H_\tau\approx F_\tau$
valid for (negative log-)likelihood objectives near the optimum~\cite{martens2015optimizing},
and restricting the displacement to the adapter parameters $\delta W$, we obtain
Eq.~\eqref{eq:bound1}.
\end{proof}

Because the classifier head is shared and small, the dominant source of forgetting is the
adapter displacement; Eq.~\eqref{eq:bound1} isolates exactly this term.

\subsection{FOLoRA minimizes the bound}
\label{sec:minimize}

Summing Eq.~\eqref{eq:bound1} over all previous tasks $\tau<t$ yields
\begin{equation}
\sum_{\tau<t}\Delta\mathcal{L}_\tau
\;\lesssim\;
\frac12\,\delta W^{\!\top}\Big(\textstyle\sum_{\tau<t} F_\tau\Big)\delta W
\;=\;
\frac12\,\delta W^{\!\top}\bar{F}\,\delta W,
\label{eq:sum}
\end{equation}
where we identify the accumulated Fisher $\sum_{\tau<t}F_\tau$ with the accumulated
importance kernel $\bar{F}$ of Eq.~\eqref{eq:acc} (restricted to the adapter parameters). Using the eigendecomposition $\bar{F}=\sum_j\sigma_j u_j u_j^{\!\top}$,
\begin{equation}
\sum_{\tau<t}\Delta\mathcal{L}_\tau
\;\lesssim\;
\frac12\sum_{j}\sigma_j\,\big(u_j^{\!\top}\delta W\big)^2.
\label{eq:eigenbound}
\end{equation}

Equation~\eqref{eq:eigenbound} shows that forgetting is bounded by a sum of terms
$\sigma_j(u_j^{\!\top}\delta W)^2$, i.e., the squared projection of the displacement onto
each eigendirection $u_j$, weighted by its importance $\sigma_j$. The FOLoRA regularizer
in Eq.~\eqref{eq:loss},
\begin{equation}
\Omega(\delta W)=\sum_j\sigma_j\,\|u_j^{\!\top}\delta W\|_F^2,
\qquad \delta W=\Delta W-\Delta W^{t-1},
\end{equation}
is precisely the Lagrangian relaxation of the constraint that the adapter displacement
remain orthogonal to the high-importance eigendirections: minimizing $\Omega$ drives
$\|u_j^{\!\top}\delta W\|$ toward zero for the largest $\sigma_j$, which in turn minimizes
the dominant terms of Eq.~\eqref{eq:eigenbound}. We therefore interpret FOLoRA as the
algorithm that \emph{directly optimizes an upper bound on catastrophic forgetting}, with
$\lambda$ controlling the stability--plasticity trade-off and $k$ controlling the rank of
the protected subspace. Following the standard single-reference approximation used by
EWC~\cite{kirkpatrick2017overcoming}, the algorithm stores one adapter snapshot
$\Delta W^{t-1}$ (after the previous task) and uses it as the reference for all past
tasks; the per-task references $\Delta W^{\tau}$ in Eq.~\eqref{eq:sum} are replaced by
$\Delta W^{t-1}$.

\subsection{Connection to equal-weight orthogonalization}
\label{sec:equalweight}

If all eigenvalues are set to $\sigma_j=1$, the bound in Eq.~\eqref{eq:eigenbound}
becomes $\frac12\|\bar{U}^{\!\top}\delta W\|^2$ with an equal-weight orthonormal basis
$\bar{U}$, which is the geometry underlying O-LoRA-style orthogonal subspace methods.
FOLoRA's eigenvalue weighting is strictly more expressive: it can concentrate protection
on the few high-curvature directions and thereby achieve a smaller forgetting for the
same rank budget, or equivalently the same forgetting with a smaller protected rank $k$.
This prediction is verified empirically by the ablation over $\lambda$ and $k$ in
Sec.~\ref{sec:ablation}.
